% TOP_PROBLEM: {"id": 30006898, "title": "Bourgain's unbounded Sobolev orbit problem for cubic defocusing NLS on the square two-torus", "category_id": 9, "category": {"id": 9, "name": "pde", "display_name": "Partial Differential Equations", "description": "PDEs and their applications in physics and geometry.", "slug": "pde", "order_index": 9, "created_at": "2026-07-31T15:26:25.671Z"}, "status": "open"}
% FIRST_POSTED: 2026-09-27
% !TeX program = lualatex
% ATTEMPT_STATUS: unresolved
% Model: GPT-6 Astra Ultra; continuation dated 27 September 2026.
\documentclass[11pt]{article}
\usepackage[margin=25mm]{geometry}
\usepackage{fontspec}
\setmainfont{Latin Modern Roman}
\usepackage{amsmath,amssymb,mathtools}
\usepackage{unicode-math}
\setmathfont{Latin Modern Math}
\usepackage{xurl,hyperref}
\hypersetup{colorlinks=true,urlcolor=blue}
\setcounter{secnumdepth}{0}
\setlength{\parindent}{0pt}
\setlength{\parskip}{5pt}
\setlength{\emergencystretch}{3em}
\begin{document}
\section*{\texorpdfstring{Bourgain's unbounded Sobolev orbit problem for cubic defocusing NLS on the square two-torus}{Bourgain's unbounded Sobolev orbit problem for cubic defocusing NLS on the square two-torus}}\label{30006898}
\subsection{Definitions and mathematical statement}
On the square torus $\mathbb T^2=({\mathbb{R}}/2\pi{\mathbb{Z}})^2$, consider the defocusing cubic equation
\[
 i\partial_tu+\Delta u=|u|^2u.
\]
For Fourier coefficients $\widehat u(k)$, take $\|u\|_{H^s}^2=\sum_{k\in{\mathbb{Z}}^2}(1+|k|^2)^s|\widehat u(k)|^2$ with a fixed harmless normalization. The question asks whether there exist one global smooth solution and one real $s>1$ such that
\[
 \limsup_{t\to+\infty}\|u(t)\|_{H^s}=+\infty.
\]
The same solution must exhibit arbitrarily large growth at arbitrarily late times. Choosing a different initial condition for each desired amplification does not establish this orbit statement. Irrational tori or equations with a convolution potential are related variants, not the fixed square-torus equation here.
\par\smallskip\textit{Formulation and concept references:} \href{https://arxiv.org/html/2308.13468}{[S1]}.

\subsection{Short English statement}
Can one smooth solution of cubic defocusing NLS on the square two-torus develop unbounded higher Sobolev norms along arbitrarily late times?
\subsection{Research attempt}
\paragraph{A single orbit, rather than a sequence of experiments.}
The inspected introduction of [S1] distinguishes the unbounded-orbit
question from arbitrary but finite growth. Its quoted finite-growth
theorem for the square torus gives a solution depending on the desired
amplification. Its own convolution-potential and irrational-torus
results are not substituted for the selected equation. We retain the
same square torus, the same cubic defocusing nonlinearity and a single
smooth solution throughout the target.

This attempt derives the conservation constraints, identifies the
rectangle resonance that permits redistribution, and states an exact
topological condition that would turn finite growth into a single
unbounded orbit. The missing step is verifying that condition, or a
uniform infinite cascade, for this equation.

\paragraph{Conserved quantities and what they actually control.}
Use normalized Haar measure on the torus so that Fourier Parseval
has no volume factor. For a smooth solution, multiplication by
$\bar u$ and integration give conservation of mass
\[
 M(u)=\int|u|^2.
\]
Indeed $u_t=i\Delta u-i|u|^2u$, and the real part of
$\int\bar u u_t$ vanishes after integration by parts. The energy is
\[
 E(u)=\frac12\int|\nabla u|^2+\frac14\int|u|^4.
 \tag{387.1}
\]
Its derivative equals
$\operatorname{Re}\int(-\Delta\bar u+|u|^2\bar u)u_t=0$
by the equation. Thus
\[
 \|u(t)\|_{H^1}^2\leq M(u(0))+2E(u(0))
 \tag{387.2}
\]
for all times on which the smooth solution exists. These identities
are compatible with the established global smooth theory used in
the formulation; they alone do not prove a uniform $H^s$ bound
for $s>1$.

The gap cannot be filled by an inequality bounding $H^s$ using only
the values on the right of (387.2). The smooth functions
$v_N(x)=N^{-1}e^{iNx_1}$ have uniformly bounded mass and energy,
but
\[
 \|v_N\|_{H^s}=N^{-1}(1+N^2)^{s/2}\longrightarrow\infty
 \quad(s>1).
 \tag{387.3}
\]
This is a statement about possible states, not a construction of
a trajectory visiting them. The equation may have dynamical
constraints beyond the two conserved scalar quantities.

Conversely, every plane wave
\[
 u(t,x)=a\exp\bigl(i k\cdot x-i(|k|^2+|a|^2)t\bigr),
 \qquad k\in\mathbb Z^2,
 \tag{387.4}
\]
is an exact global solution and has constant Sobolev norms of every
order. Thus a high-frequency initial state by itself produces no
growth. The problem requires transfer among frequencies over time.

\paragraph{No fixed Fourier truncation can display the target behavior.}
Suppose a solution of any mass-conserving finite-mode model remains
supported in $|k|\leq K$. Then
\[
 \|u(t)\|_{H^s}^2
 \leq(1+K^2)^{s-1}\|u(t)\|_{H^1}^2
 \quad(s>1).
 \tag{387.5}
\]
In particular the Hamiltonian Galerkin truncation, obtained by
orthogonally projecting the nonlinearity to these modes, conserves
the restrictions of mass and energy. To check energy conservation,
its vector field is the Hamiltonian gradient of (387.1) on the
finite Fourier subspace, so the same calculation works with the
orthogonal projection inserted. The mass conservation is equally
direct. Its finite-dimensional trajectories therefore have a uniform
$H^s$ bound depending on $K$ and the initial invariants.

Finite-mode computations can detect amplification and test a local
transfer mechanism. They cannot prove an unbounded norm on a fixed
orbit unless there is an argument passing to infinitely many active
scales while maintaining a single admissible initial condition.
Letting $K$ increase between experiments changes the system being
tested and does not remove this requirement.

\paragraph{The exact resonant geometry on the square torus.}
Write $u(t,x)=\sum_k u_k(t)e^{ik\cdot x}$ and remove the linear
oscillation by $a_k(t)=e^{i|k|^2t}u_k(t)$. The Fourier equation is
\[
 i\dot a_k=\sum_{k_1-k_2+k_3=k}
 a_{k_1}\bar a_{k_2}a_{k_3}
 e^{-it\Omega},\qquad
 \Omega=|k_1|^2-|k_2|^2+|k_3|^2-|k|^2.
 \tag{387.6}
\]
Using the momentum relation to expand the last square gives
\[
 \Omega=-2(k_1-k_2)\cdot(k_3-k_2).
 \tag{387.7}
\]
Thus a nondegenerate resonant quadruple consists of the vertices
of a lattice rectangle. For example
\[
 0,\quad a=(N,0),\quad b=(0,N),\quad d=a+b
\]
obey $a\cdot b=0$ and
$|0|^2+|d|^2=|a|^2+|b|^2$. This equality permits an exchange
between the opposite pairs without changing the quadratic energy.
For the Sobolev weight $w(k)=(1+|k|^2)^s$, strict convexity gives
\[
 w(0)+w(d)-w(a)-w(b)
       =1+(1+2N^2)^s-2(1+N^2)^s>0
       \quad(s>1).
 \tag{387.8}
\]
Thus the higher norm is not algebraically forced to remain constant
by this particular resonance.

To make that observation dynamical without overstating it, consider
only the four-mode interaction Hamiltonian
\[
 H_{\rm int}=Z+\bar Z,\qquad
 Z=a_0a_d\bar a_a\bar a_b,
 \qquad i\dot a_j=\partial_{\bar a_j}H_{\rm int}.
\]
Writing $I_j=|a_j|^2$ gives the exact identities
\[
 \dot I_0=\dot I_d=-2\operatorname{Im}Z,\qquad
 \dot I_a=\dot I_b=2\operatorname{Im}Z.
 \tag{387.9}
\]
They follow from
$\dot I_j=2\operatorname{Im}(\bar a_j\partial_{\bar a_j}H_{\rm int})$.
Mass and the quadratic Fourier energy are preserved, while the
weighted sum has derivative
\[
 \frac{d}{dt}\sum_jw(j)I_j
 =2\operatorname{Im}Z\,[w(a)+w(b)-w(0)-w(d)].
 \tag{387.10}
\]
Suitable phases make this nonzero. This is a verified local mechanism
in the stated interaction model. It is not asserted that the full
NLS stays on this four-mode subsystem: other cubic interactions and
nonresonant terms are present in (387.6).

\paragraph{Why averaging has a finite-time qualification.}
Terms with $\Omega\ne0$ in (387.6) oscillate. Integrating one such
term by parts produces a factor $1/\Omega$, a boundary term, and
an integral involving derivatives of the amplitudes. Even though
the square-torus frequencies make nonzero $\Omega$ integers, the
number of interacting terms, amplitude sizes and duration all enter
the remainder estimate. Dropping oscillatory terms is therefore a
finite-time approximation requiring explicit control, not an exact
identity for arbitrarily long trajectories.

Suppose a resonant construction executes $J$ transfers and is
shadowed by the full equation with an error bound depending on $J$.
To take $J\to\infty$, one needs compatible initial data, summable
errors and existence on every corresponding time interval. A theorem
giving one independently chosen shadowing solution for each finite
$J$ does not furnish these properties. Conservation laws may help
control the low norms, but they do not by themselves guarantee
high-frequency errors remain small in the norm in which the transfer
is measured.

\paragraph{The precise logical gap in finite amplification.}
Finite growth has the quantifier pattern
\[
 \text{for every }R\text{ there exist }u_R\text{ and }T_R
       \text{ with }\|u_R(T_R)\|_{H^s}>R.
\]
The question instead asks for one $u$ for which the inequality holds
for every $R$ at some late time. Interchanging these quantifiers is
not valid. The elementary family of functions
$F_j(t)=j\sin^2t$ has arbitrarily large values across the family,
while each individual function is bounded. It is a logical
countercheck, not an NLS model.

Compactness of a sequence of initial data is insufficient as well.
For example, a finite-growth theorem may choose
$\|u_j(0)\|_{H^s}\to0$ while $\|u_j(T_j)\|_{H^s}\to\infty$.
Continuity of the solution map on each fixed bounded time interval
then makes the limiting solution the zero solution on every such
interval. Necessarily the growth times escape those intervals.
The large values at $T_j\to\infty$ do not pass through convergence
that is only uniform on compact time intervals.

\paragraph{An exact Baire-category route, with the missing premise exposed.}
Let $\mathcal X=C^\infty(\mathbb T^2;\mathbb C)$ with its usual
complete metrizable Frechet topology. Use the established global
smooth flow and its continuity on fixed time intervals. Fix $s>1$
and define
\[
 \mathcal U_m=\bigcup_{t\geq0}
      \{u_0\in\mathcal X:\|u(t;u_0)\|_{H^s}>m\}.
 \tag{387.11}
\]
Each $\mathcal U_m$ is open, since it is a union of inverse images
of open norm intervals under continuous fixed-time solution maps.
If every $\mathcal U_m$ were dense, Baire's theorem would give a
dense intersection $\bigcap_{m\geq1}\mathcal U_m$. Every point in
that intersection would generate a smooth orbit unbounded in $H^s$.
Continuity of one smooth orbit bounds its norm on each compact time
interval, so its arbitrarily large values must occur at arbitrarily
late times; the required limsup would follow.

This reduction cleanly specifies the needed strengthening:
large growth must be available in \emph{every} nonempty smooth
neighborhood, not merely from some small data near the origin.
Such a neighborhood can prescribe closeness in an arbitrarily high
but finite list of Sobolev norms. A high-frequency perturbation small
in one chosen $H^s$ need not be small in that neighborhood. The
finite-growth theorem cited in [S1] does not assert this density
premise, and the present calculations do not establish it.

An alternative constructive route would select nested nonempty
closed neighborhoods of initial data, each guaranteeing a new
transfer at a specified finite time, with diameters tending to zero
in a complete smooth metric. Their intersection would provide one
orbit. Again the hard requirement is the ability to perform the next
transfer while preserving all previous finite-time requirements and
remaining in the selected neighborhood. Stating nestedness without
constructing it would simply restate the desired infinite cascade.

\paragraph{Diagnostics and outcome.}
The finite checks verify the rectangle identity (387.7), the
strict higher-weight imbalance (387.8) for sample exponents, and
the mass and quadratic-energy cancellations in (387.9). The
fixed-cutoff inequality (387.5) is proved above rather than tested
by the script. These are finite arithmetic diagnostics; no finite
Fourier simulation is reported as an
unbounded-orbit construction.

The conservation constraints permit higher-norm growth, and a
resonant interaction can change the required weighted energy.
The first unresolved step is turning such local or finite transfer
information into the uniform density or compatible infinite
shadowing required for one smooth square-torus solution. The
selected problem remains unresolved by this attempt.
\subsection{Sources}
\begin{itemize}
\item[S1]Filippo Giuliani, Sobolev instability in the cubic NLS equation with convolution potentials on irrational tori (2023 preprint; JDE 2025).
\url{https://arxiv.org/html/2308.13468}.
\textit{Formulation source.}
\end{itemize}
\end{document}
